\lecture{3}{Fri 17 Feb 2023 16:38}{Outer measure of interval}

\begin{definition}[Lebesgue outer measure]
	If $E \subset \mathbb{R}$ we define 
	\[
		m^*(E) = \operatorname{inf} \{\sum_{k=1}^{\infty} |I_k| : I_k \text{ interval }, \bigcup_{k=1} ^\infty I_k \supset E\} 
	.\] 
\end{definition}

\begin{proposition}
	If $I \subset \mathbb{R}$ is an interval then $m^*(I) = |I|$.
\end{proposition}
\begin{proof}
	We show for $I = [0,1]$. We will show if $\bigcup_{k \in \mathbb{N}} I_k \supset [0,1]$ then $\sum_{k=1}^{\infty} |I_k| \ge 1$. Indeed: since $[0,1]$ compact, we can choose finite $K \subset \mathbb{N}$ such that $\bigcup_{k \in K} I_k \supset [0,1]$. Let $k_1 \in K$ be such that $0 \in I_{k_1} = (a_1, b_1)$.
	
	Now if $b_{k_i} \le 1$, let $k_{i+1}$ be such that $b_{k_i}\in I_{k_{i+1}} = \left( a_{k_{i+1}}, b_{k_{i+1}} \right) $.
	Such process must terminate since 1) we always get a different interval, 2) $K$ is finite. Now
	\begin{align*}
		\sum_{k=1}^{\infty} |I_k| &\ge 
		\sum_{i=1}^{n} \left| I_{k_i} \right| \\
		&= (b_1-a_1) + \ldots + (b_n - a_n) \\
		&\ge (b_1-0) + (b_2-b_1) + \ldots + (b_n - b_{n-1}) \\
		&= b_n >1 \\
	.\end{align*} 
	Hence $m^*([0,1]) \ge 1$. Also $m^*([0,1]) \le 1$, as
	\[
		I_1 = [0,1], I_{k} = \emptyset , k\ge 2
	.\] 
	is a cover of $I$.
\end{proof}

\begin{proposition}(Countable subadditivity of Lebesgue inner measure)
	\[
		m^*\left(\bigcup_{k \in \mathbb{N}} E_k\right) \le  \sum_{k=1}^{\infty} m^*(E_k)
	.\] 
\end{proposition}
\begin{proof}
	Let $\varepsilon>0$. For each $E_k$ we cover it by intervals $I_j^k$ such that $\sum_{j=1}^{\infty} \left| I_j^k \right| < m^*(E_k) + \frac{\varepsilon}{2^k}$.

	Then $\bigcup_{j,k \in \mathbb{N}} I_j^k \supset \bigcup_{k \in \mathbb{N}} E_k$, so
	\begin{align*}
		m^*(\bigcup_{k \in \mathbb{N}} E_k)
		&\le \sum_{j=1}^{\infty} \sum_{k=1}^{\infty} \left| I_j^k \right|  \\
		&= \sum_{k=1}^{\infty} \sum_{j=1}^{\infty} \left| I_j^k \right|  \\
		&< \sum_{k=1}^{\infty} \left( m^*(E_k) + \frac{\varepsilon}{2^k} \right)  \\
		&= \sum_{k=1}^{\infty} m^*(E_k) + \varepsilon
	.\end{align*}
\end{proof}

\begin{definition}[Lebesgue measurability]
	Let $E \subset [0,1]$. We define
	\[
		m_*(E) = 1 - m^*([0,1] \setminus E) \le  m^*(E)
	.\] 
	We say $E$ is \textit{measurable} if $m^*(E) = m_*(E)$.
\end{definition}
c.f. [[@taoIntroductionMeasureTheory2011]] page 17, 28

\begin{definition}[Carathéodory's definition of measurability]
	We say $E \subset \mathbb{R}$ is measurable if for any $A \subset \mathbb{R}$, 
	\[
		m^*(A) = m^*(A \cap E) + m^*(A \setminus E)
	.\] 
\end{definition}
c.f. page 33
\begin{intuition}
	$E$ cuts $A$ clearly.
\end{intuition}

\begin{remark}
	$\le $ is automatic; so it suffices to require $\ge $. Thus it suffices to require $m^*(A)$ finite.
\end{remark}

\begin{example}
	\begin{itemize}
		\item If $m^*(E) = 0$ then $E$ measurable.
	\end{itemize}
\end{example}

